\begin{frame}{왜 $\tau = t - n + 1$인가 --- 갱신 시차}
  \begin{itemize}
    \item $\tau$(갱신 대상 시점)는 항상 현재 $t$보다 \textbf{$n-1$만큼
      뒤처짐} --- $n$ 스텝의 실제 보상이 다 모여야 목표값을 계산
      가능.
    \item $t$가 한 스텝씩 가면 갱신 차례인 $\tau$도 한 칸씩 앞으로
      밀려남.
    \item 에피소드 종료(``$t \ge T$'') 이후에도 이 밀림이
      \texttt{tau == T-1}까지 계속됨 --- 아직 $n$ 스텝 목표값을
      계산 못 한 상태가 남아 있어서(꼬리 갱신).
  \end{itemize}
  \vskip0.8em
  \begin{block}{갱신 횟수는 변하지 않는다}
    에피소드가 $T$ 이동 상태면 $n=1$이든 $n=T+1$이든 \textbf{모든
    이동 상태마다 정확히 1번씩}(총 $T$번) 갱신된다. $n$이 바꾸는
    것은 갱신 \textbf{횟수}가 아니라, 첫 갱신까지의 \textbf{시차}와
    목표값의 \textbf{구성}(실제/추정 비율).
  \end{block}
\end{frame}
